\documentclass[12pt,a4paper]{article}

\usepackage[utf8]{inputenc}
\usepackage[T1]{fontenc}
\usepackage[portuguese]{babel}
\usepackage{geometry}
\usepackage{graphicx}
\usepackage{booktabs}
\usepackage{array}
\usepackage{xcolor}
\usepackage{listings}
\usepackage{hyperref}
\usepackage{amsmath}
\usepackage{enumitem}
\usepackage{titlesec}
\usepackage{fancyhdr}
\usepackage{microtype}

\geometry{top=2.5cm, bottom=2.5cm, left=2.8cm, right=2.8cm}

% ---------- cores e estilo de código ----------
\definecolor{codebg}{RGB}{245,245,245}
\definecolor{codeframe}{RGB}{200,200,200}
\definecolor{keyword}{RGB}{0,0,180}
\definecolor{comment}{RGB}{80,120,80}
\definecolor{string}{RGB}{160,30,30}
\definecolor{darkblue}{RGB}{0,40,100}

\lstset{
  backgroundcolor=\color{codebg},
  frame=single,
  rulecolor=\color{codeframe},
  basicstyle=\ttfamily\small,
  keywordstyle=\color{keyword}\bfseries,
  commentstyle=\color{comment}\itshape,
  stringstyle=\color{string},
  breaklines=true,
  columns=fullflexible,
  keepspaces=true,
  showstringspaces=false,
  numbers=none,
}

% ---------- cabeçalho ----------
\pagestyle{fancy}
\fancyhf{}
\rhead{\small AED 2025--26}
\lhead{\small Grafo de Conhecimento CTI}
\cfoot{\thepage}
\renewcommand{\headrulewidth}{0.4pt}

% ---------- título ----------
\title{%
  \vspace{-1cm}
  {\large Escola Naval \textemdash\ Curso EN-AEL}\\[0.4cm]
  {\large Algoritmos e Estruturas de Dados 2025--26}\\[0.8cm]
  {\LARGE\bfseries Grafo de Conhecimento para\\[0.2cm]
   Cyber Threat Intelligence}\\[0.4cm]
  {\large Relatório Técnico \textemdash\ Nível 3: Algoritmos de Análise}
}
\author{Daniel Zayika}
\date{\today}

\begin{document}

\maketitle
\thispagestyle{fancy}

\tableofcontents
\newpage

% ============================================================
\section{Introdução}

Este relatório documenta a implementação e os resultados do \textbf{Nível~3} do projeto, que acrescenta algoritmos de análise de grafos ao sistema de Knowledge Graph para Cyber Threat Intelligence (CTI) desenvolvido nos níveis anteriores.

O sistema representa entidades de segurança (atores, malware, indicadores, campanhas, vulnerabilidades) como nós de um grafo dirigido, com relações tipadas entre elas. O Nível~3 transforma este repositório de dados numa \emph{ferramenta analítica}: é agora possível atravessar o grafo, descobrir caminhos entre entidades, quantificar a relevância de cada nó e resolver aliases de atores reportados sob nomes diferentes.

O dataset sintético utilizado contém \textbf{22 nós} e \textbf{26 relações}, modelando dois cenários reais: o ataque \emph{SolarWinds Supply Chain} (2020), atribuído ao APT29/Cozy Bear, e a campanha \emph{WannaCry} (2017), atribuída ao Lazarus Group.

% ============================================================
\section{Estrutura do Código}

O código está organizado nos seguintes módulos:

\begin{center}
\begin{tabular}{lp{9cm}}
\toprule
\textbf{Ficheiro} & \textbf{Responsabilidade} \\
\midrule
\texttt{src/types.h}           & Enumerações \texttt{EntityType} e \texttt{RelationType}; matriz de relações válidas \\
\texttt{src/node.h}            & Estrutura \texttt{Node} (id, label, tipo, propriedades) \\
\texttt{src/edge.h}            & Estrutura \texttt{Edge} (origem, destino, relação, descrição) \\
\texttt{src/graph.h/.cpp}      & Classe \texttt{Graph}: lista de adjacência, inserção, consultas \\
\texttt{src/database.h/.cpp}   & Camada SQLite: persistência e carregamento do grafo \\
\texttt{src/algorithms.h/.cpp} & Algoritmos: BFS, DFS, caminho mais curto, PageRank, resolução de entidades, consultas analíticas \\
\texttt{src/synthetic\_data.h/.cpp} & Dataset sintético com entidades e relações CTI reais \\
\texttt{src/tests.h/.cpp}      & Suite de testes automáticos \\
\texttt{main.cpp}              & Ponto de entrada; demonstração em consola \\
\bottomrule
\end{tabular}
\end{center}

% ============================================================
\section{Algoritmos Implementados}

\subsection{Travessia em Largura (BFS)}

A BFS (\emph{Breadth-First Search}) explora o grafo nível a nível a partir de um nó inicial, utilizando uma fila (\texttt{std::queue}). Visita primeiro todos os vizinhos diretos, depois os vizinhos dos vizinhos, e assim sucessivamente.

\medskip
\noindent\textbf{Complexidade:} $O(V + E)$, onde $V$ é o número de nós e $E$ o número de arestas.

\medskip
\noindent\textbf{Resultado obtido} (partindo de \texttt{apt29}, grafo dirigido):

\begin{lstlisting}
apt29 -> camp-solarwinds -> id-nato -> id-usgov -> mimikatz -> sunburst
      -> atk-cred-dump -> atk-supply-chain -> teardrop
\end{lstlisting}

A BFS garante que os nós mais próximos de \texttt{apt29} são visitados primeiro: a campanha e os alvos diretos aparecem antes do malware e das técnicas de ataque, que se encontram a maior distância no grafo.

\subsection{Travessia em Profundidade (DFS)}

A DFS (\emph{Depth-First Search}) explora cada ramo do grafo até ao fim antes de recuar, utilizando uma pilha (\texttt{std::stack}).

\medskip
\noindent\textbf{Complexidade:} $O(V + E)$.

\medskip
\noindent\textbf{Resultado obtido} (partindo de \texttt{apt29}, grafo dirigido):

\begin{lstlisting}
apt29 -> camp-solarwinds -> id-usgov -> sunburst -> atk-supply-chain
      -> teardrop -> id-nato -> mimikatz -> atk-cred-dump
\end{lstlisting}

Ao contrário da BFS, a DFS segue imediatamente a cadeia \texttt{camp-solarwinds} $\to$ \texttt{sunburst} $\to$ \texttt{atk-supply-chain} $\to$ \texttt{teardrop} em profundidade, revelando a cadeia de ataque completa antes de voltar para os alvos.

\subsection{Caminho Mais Curto}

Implementado com BFS sobre o grafo não ponderado (cada aresta tem custo unitário). O algoritmo regista o nó pai de cada nó visitado e reconstrói o caminho por retrocesso (\emph{backtracking}).

\medskip
\noindent\textbf{Complexidade:} $O(V + E)$.

\medskip
\noindent\textbf{Resultados obtidos:}

\begin{center}
\begin{tabular}{lll}
\toprule
\textbf{Origem} & \textbf{Destino} & \textbf{Caminho} \\
\midrule
\texttt{ioc-hash-sunburst} & \texttt{apt29}       & \texttt{ioc-hash-sunburst} $\to$ \texttt{sunburst} $\to$ \texttt{apt29} \\
\texttt{ioc-ip-91}         & \texttt{id-nato}     & \texttt{ioc-ip-91} $\to$ \texttt{apt28} $\to$ \texttt{id-nato} \\
\texttt{ioc-hash-wannacry} & \texttt{vuln-ms17-010} & \texttt{ioc-hash-wannacry} $\to$ \texttt{wannacry} $\to$ \texttt{vuln-ms17-010} \\
\bottomrule
\end{tabular}
\end{center}

\medskip
Os três caminhos têm comprimento~2 (três nós), o que reflete a estrutura do dataset: os indicadores ligam-se diretamente ao malware/ator, que por sua vez se liga ao alvo ou vulnerabilidade. Estes caminhos curtos são exatamente o que um analista CTI precisa para justificar uma atribuição.

\subsection{PageRank}

O PageRank quantifica a importância de cada nó com base na estrutura de ligações do grafo. Um nó é relevante se muitos nós relevantes apontam para ele. A fórmula iterativa é:

\[
  PR(u) = \frac{1-d}{N} + d \sum_{v \in \text{In}(u)} \frac{PR(v)}{\text{Out}(v)}
\]

onde $d = 0{,}85$ é o fator de amortecimento, $N$ o número total de nós, $\text{In}(u)$ o conjunto de nós que apontam para $u$ e $\text{Out}(v)$ o grau de saída de $v$. O algoritmo inclui tratamento de \emph{dangling nodes} (nós sem arestas de saída), redistribuindo a sua massa uniformemente.

\medskip
\noindent\textbf{Complexidade por iteração:} $O(V + E)$. Converge tipicamente em menos de 100 iterações.

\medskip
\noindent\textbf{Top-8 entidades por PageRank:}

\begin{center}
\begin{tabular}{clll}
\toprule
\textbf{\#} & \textbf{Entidade} & \textbf{Score} & \textbf{Tipo} \\
\midrule
1 & \texttt{sunburst}         & 0.114667 & malware \\
2 & \texttt{camp-solarwinds}  & 0.096785 & campaign \\
3 & \texttt{atk-supply-chain} & 0.070113 & attack-pattern \\
4 & \texttt{teardrop}         & 0.070113 & malware \\
5 & \texttt{vuln-ms17-010}    & 0.070058 & vulnerability \\
6 & \texttt{id-usgov}         & 0.069236 & identity \\
7 & \texttt{wannacry}         & 0.057269 & malware \\
8 & \texttt{atk-cred-dump}    & 0.045267 & attack-pattern \\
\bottomrule
\end{tabular}
\end{center}

\medskip
\noindent\textbf{Interpretação:} O \texttt{SUNBURST} lidera com grande margem porque é referenciado por múltiplas entidades de alto peso: \texttt{apt29}, \texttt{dukes}, \texttt{camp-solarwinds} e dois indicadores IoC. A campanha \emph{SolarWinds} ocupa o segundo lugar por ser o hub central do maior cluster do grafo. A \texttt{vuln-ms17-010} (EternalBlue) surge no top-5 apesar de pertencer ao cluster WannaCry, que é menor, porque é um ponto de convergência único nesse cluster — apenas o WannaCry aponta para ela, mas o WannaCry tem alto score por si próprio.

\subsection{Resolução de Entidades}

O sistema implementa resolução de aliases com dois mecanismos:

\begin{enumerate}
  \item \textbf{\texttt{resolveByName}}: normaliza a query (minúsculas, apenas alfanuméricos) e compara com todas as variantes de cada nó: id, label, partes do label separadas por~``/'', e a propriedade \texttt{alias}.

  \item \textbf{\texttt{findPotentialDuplicates}}: agrupa nós que partilham pelo menos uma variante normalizada em comum.
\end{enumerate}

\medskip
\noindent\textbf{Resultados obtidos:}

\begin{lstlisting}
resolveByName("APT29")     -> [apt29] "APT29 / Cozy Bear"
resolveByName("Cozy Bear") -> [apt29] "APT29 / Cozy Bear"
                              [dukes] "The Dukes"

Duplicados detectados: { apt29, dukes }
\end{lstlisting}

O nó \texttt{dukes} foi introduzido propositadamente no dataset para simular uma segunda fonte de inteligência que reporta o mesmo ator sob o nome \emph{The Dukes}, com \texttt{alias = "Cozy Bear"}. O algoritmo deteta corretamente que ambos partilham o alias e os agrupa como potenciais duplicados — situação frequente em CTI real, onde diferentes feeds usam nomenclaturas distintas para o mesmo ator.

% ============================================================
\section{Consultas Analíticas}

\subsection{Atores que usam determinado malware}

\begin{center}
\begin{tabular}{ll}
\toprule
\textbf{Malware} & \textbf{Atores} \\
\midrule
\texttt{sunburst}  & \texttt{apt29} (APT29/Cozy Bear), \texttt{dukes} (The Dukes) \\
\texttt{wannacry}  & \texttt{lazarus} (Lazarus Group) \\
\bottomrule
\end{tabular}
\end{center}

A consulta para \texttt{sunburst} devolve dois atores porque o dataset contém intencionalmente o nó duplicado \texttt{dukes}. Após resolução de entidades, conclui-se que ambos representam o mesmo ator.

\subsection{Indicadores que apontam para determinada ameaça}

\begin{center}
\begin{tabular}{ll}
\toprule
\textbf{Ameaça} & \textbf{Indicadores (IoCs)} \\
\midrule
\texttt{sunburst} & \texttt{ioc-domain-avsvmcloud} (\texttt{avsvmcloud[.]com}) \\
                  & \texttt{ioc-hash-sunburst} (SHA-256 da DLL maliciosa) \\
\bottomrule
\end{tabular}
\end{center}

Estes dois IoCs são suficientes para detetar a presença do SUNBURST numa rede: o hash identifica o binário e o domínio identifica o servidor C2.

% ============================================================
\section{Cenário de Investigação}

O sistema foi demonstrado com o seguinte cenário narrativo, típico de uma investigação CTI:

\begin{enumerate}
  \item Um analista deteta nos logs o hash \texttt{32519b85...} (indicador \texttt{ioc-hash-sunburst}).
  \item O sistema resolve: o indicador \emph{indica} o malware \texttt{SUNBURST}.
  \item Consulta analítica: que atores usam o SUNBURST? $\to$ \texttt{apt29} e \texttt{dukes}.
  \item Resolução de entidades: \texttt{apt29} e \texttt{dukes} são o mesmo ator (APT29/Cozy Bear/The Dukes).
  \item Caminho mais curto que justifica a atribuição:
  \[
    \texttt{ioc-hash-sunburst} \to \texttt{sunburst} \to \texttt{apt29}
  \]
  \item Conclusão: o incidente é atribuível ao APT29, grupo de espionagem russo que comprometeu agências do governo norte-americano através da supply chain da SolarWinds.
\end{enumerate}

Este cenário demonstra como o grafo transforma dados dispersos (um hash isolado) em conhecimento estruturado e acionável (atribuição a um ator com justificação pelo caminho de relações).

% ============================================================
\section{Testes}

A suite de testes automáticos (\texttt{./cti\_graph --testes}) cobre:

\begin{itemize}
  \item Inserção e validação de nós e relações
  \item Vizinhos e filtragem por tipo de entidade
  \item Integridade do dataset sintético
  \item Persistência SQLite (round-trip guardar/carregar)
  \item BFS, DFS, caminho mais curto, PageRank
  \item Resolução de entidades e deteção de duplicados
  \item Consultas analíticas
\end{itemize}

\medskip
\noindent Todos os testes passaram com sucesso na execução documentada neste relatório.

% ============================================================
\section{Conclusões}

O Nível~3 transforma o sistema de um repositório passivo num motor de análise. Os principais resultados são:

\begin{itemize}
  \item O \textbf{SUNBURST} é a entidade mais central do grafo (PageRank = 0.1147), refletindo o papel pivô que desempenhou no ataque SolarWinds como ponto de convergência entre o ator, a campanha e os indicadores.

  \item O \textbf{caminho mais curto} permite justificar atribuições de forma auditável: em apenas dois saltos é possível ligar um IoC ao ator responsável.

  \item A \textbf{resolução de entidades} é essencial em CTI real, onde o mesmo ator é frequentemente reportado sob nomes diferentes por fontes distintas. O mecanismo implementado deteta corretamente o alias \emph{Cozy Bear} como denominador comum entre \texttt{apt29} e \texttt{dukes}.

  \item A combinação de BFS/DFS com consultas analíticas permite responder às perguntas fundamentais de threat intelligence: \emph{quem usa o quê}, \emph{o que aponta para quê}, e \emph{qual o caminho entre dois pontos do grafo}.
\end{itemize}

\end{document}
